\chapter{Grammar, Syntax, and Examples of JSONPath}

\section{Formal Grammar}

JSONPath expressions are based on a formal grammar that provide the syntax rules for querying JSON data. The grammar make sure that JSONPath expressions are parsable by JSON parsers.  
A simple version of the JSONPath grammar, inspired by the IETF drafts \cite{ietf_jsonpath_draft, ietf_rfc9535} and implementations like PostgreSQL’s \texttt{jsonpath\_gram.y} \cite{postgres_jsonpath}, is given below using a Backus-Naur Form (BNF)-style notation:

\begin{quote}
	\begin{tabbing}
		jsonpath       ::= root-selector selector-list \\
		root-selector  ::= '\$' 
		selector-list  ::= selector selector-list \textbar\ $\epsilon$ \\
		selector       ::= '.' child-name \\
		\quad \textbar\ '..' recursive-name \\
		\quad \textbar\ '[' array-index ']' \\
		\quad \textbar\ '[' array-slice ']' \\
		\quad \textbar\ '[' filter ']' \\
		child-name     ::= NAME \\
		recursive-name ::= NAME \\
		array-index    ::= INTEGER \\
		array-slice    ::= INTEGER? ':' INTEGER? (':' INTEGER?)? \\
		filter         ::= '?(' expression ')' \\
		expression     ::= \dots
	\end{tabbing}
\end{quote}

Here:  
- \texttt{\$} is the root object of the JSON.  
- \texttt{@} is the current node in filter expressions.  
- Selectors allow navigation to child elements, recursive descent, array elements, or filtered subsets.

\section{Core Syntax and Expressions}

The primary syntax components of JSONPath are available in the IETF JSONPath standard \cite{ietf_rfc9535} and popular tutorials \cite{toolsqa_jsonpath, azure_jsonpath}:

\begin{itemize}
	\item \textbf{Root Selector (\$)}: the root of the JSON data.
	\item \textbf{Current Node (@)}: the current context node, primarily used in filter expressions.
	\item \textbf{Dot Notation (.)}: Selects child members by name, e.g., \texttt{.store}.
	\item \textbf{Recursive Descent (..)}: Searches recursively for matching child elements at any depth.
	\item \textbf{Wildcards (*)}: Matches all elements in an array or all keys in an object.
	\item \textbf{Array Indexing and Slicing}:  
	- Single index: \texttt{[0]}, \texttt{[1]} selects elements by position.  
	- Slice: \texttt{[start:end:step]} selects a range of elements.
	\item \textbf{Filters ([?()])}: Allows selection based on expressions, such as comparisons or logical operations inside brackets.  
	\item \textbf{Expressions inside Filters}:  
	- Comparison operators: \texttt{==, !=, <, >, <=, >=}  
	- Logical operators: \texttt{\&\& (and), || (or), ! (not)}  
	- Functions: \texttt{length()}, \texttt{sum()}, etc.
\end{itemize}

\section{Examples}

The following JSON data represent a bookstore \cite{toolsqa_jsonpath}:

\begin{verbatim}
	{
		"store": {
			"book": [
			{
				"category": "reference",
				"author": "Nigel Rees",
				"title": "Sayings of the Century",
				"price": 8.95
			},
			{
				"category": "fiction",
				"author": "Evelyn Waugh",
				"title": "Sword of Honour",
				"price": 12.99
			},
			{
				"category": "fiction",
				"author": "Herman Melville",
				"title": "Moby Dick",
				"price": 8.99
			}
			],
			"bicycle": {
				"color": "red",
				"price": 19.95
			}
		}
	}
\end{verbatim}

Some JSONPath queries and their explanations:

\begin{itemize}
	\item \texttt{\$.store.book[*].author}  
	Selects the authors of all books in the store.  
	\textbf{Result:} \texttt{["Nigel Rees", "Evelyn Waugh", "Herman Melville"]}
	
	\item \texttt{\$.store.book[0]}  
	Selects the first book in the array.  
	\textbf{Result:} The entire JSON object for the first book.
	
	\item \texttt{\$.store.book[?(@.price < 10)].title}  
	Selects titles of books with a price less than 10.  
	\textbf{Result:} \texttt{["Sayings of the Century", "Moby Dick"]}
	
	\item \texttt{\$.store..price}  
	Recursively selects all prices in the store (books and bicycle).  
	\textbf{Result:} \texttt{[8.95, 12.99, 8.99, 19.95]}
	
	\item \texttt{\$..author}  
	Selects all authors in the document regardless of depth.  
	\textbf{Result:} \texttt{["Nigel Rees", "Evelyn Waugh", "Herman Melville"]}
\end{itemize}

These examples shows how JSONPath expressions gives powerful and flexible ways to navigate and extract data from complex JSON structures.
